\section{光滑映射}
\label{sec:mmaps-smooth-maps}

设$M$是$m$维光滑流形，$N$是$n$维光滑流形，$F:M\to N$是映射。判断$F$是否光滑，仍要借助坐标，但现在两侧都要用图：输入的点用$M$的图记录，输出的点用$N$的图记录。这样得到的是欧氏开集之间的映射，可以用多元微积分检查。

\paragraph{坐标表示与光滑性}
设$(U,\varphi)$是$M$的图，$(V,\psi)$是$N$的图，并且$F(U)\subseteq V$。把输入坐标先送回点，作用$F$，再读出输出坐标，得到
\begin{equation}
\hat F=\psi\circ F\circ\varphi^{-1}:\varphi(U)\to\psi(V)\subseteq\mathbb R^n ,
\label{eq:mmaps-coordinate-representation}
\end{equation}
称为$F$在这对图下的\term{坐标表示}（coordinate representation）。条件$F(U)\subseteq V$保证输出点能用同一张图$\psi$读出。

\begin{definition}[光滑映射]
\label{def:mmaps-smooth-map}
映射$F:M\to N$称为\term{光滑映射}（smooth map），若$F$连续，并且对每个$p\in M$，存在包含$p$的图$(U,\varphi)$与包含$F(p)$的图$(V,\psi)$，使$F(U)\subseteq V$且坐标表示\eqref{eq:mmaps-coordinate-representation}光滑。$M$或$N$带边界时，光滑性按第~\ref{chap:manifold-basics}章的约定理解：半空间开子集上的映射光滑，指它在每点附近能光滑延拓到一个欧氏开邻域上。
\end{definition}

连续性的作用是保证满足$F(U)\subseteq V$的图存在：任给包含$F(p)$的图$(V,\psi)$，$F^{-1}(V)$是包含$p$的开集，把$M$的图限制到$U\cap F^{-1}(V)$上即可。

\paragraph{与图的选择无关}
定义只要求每点附近存在一对图使坐标表示光滑。换用别的图对，结论不变：若$(U',\varphi')$与$(V',\psi')$是另一对满足条件的图，在重叠处
\[
\psi'\circ F\circ\varphi'^{-1}=(\psi'\circ\psi^{-1})\circ\hat F\circ(\varphi\circ\varphi'^{-1}),
\]
右端是两个光滑坐标变换与$\hat F$的复合，仍光滑。同样的论证说明光滑映射的复合仍是光滑映射：若$F:M\to N$与$G:N\to P$都光滑，则在合适的图下$\widehat{G\circ F}=\hat G\circ\hat F$，光滑映射的复合光滑。

\paragraph{特例：曲线与函数}
第~\ref{chap:manifold-basics}章的两类对象都是光滑映射的特例。把区间$I$看作一维流形（带端点时为带边界流形），用恒等映射作图，光滑曲线$\gamma:I\to M$就是$I$到$M$的光滑映射，坐标表示是坐标曲线$\varphi\circ\gamma$。把$\mathbb R$看作一维流形，光滑函数$f:M\to\mathbb R$就是$M$到$\mathbb R$的光滑映射，坐标表示是$\hat f=f\circ\varphi^{-1}$。今后关于一般光滑映射的结论，都可以用这两类特例检验。

\paragraph{例子：球面的旋转}
设$\bm R(\theta)$是平面上转角为$\theta$的旋转矩阵。绕竖直轴转$\theta$的映射
\[
R_\theta:S^2\to S^2,\qquad R_\theta(\bm y)=(\bm R(\theta)(y_1,y_2)^\top,\;y_3)
\]
把球面送到自身。在北极图中，$\sigma_N(R_\theta(\bm y))=\bm R(\theta)(y_1,y_2)^\top/(1-y_3)=\bm R(\theta)\sigma_N(\bm y)$，所以坐标表示是线性映射$\hat R_\theta(\bm u)=\bm R(\theta)\bm u$，光滑。北极是$R_\theta$的不动点，那里改用南极图，坐标表示同样是$\bm w\mapsto\bm R(\theta)\bm w$。于是$R_\theta$是光滑映射。

绕水平轴的旋转会把北极移走，这时需要按点选择图。例如绕$y_1$轴转$\pi/2$的映射$Q(\bm y)=(y_1,-y_3,y_2)$把北极送到$(0,-1,0)$，把点$(0,1,0)$送到北极。在$(0,1,0)$附近，输入可以用北极图，输出却必须用南极图，因为北极不在北极图的定义域内。定义中“对每个$p$存在一对图”正是为这种情况准备的。

\paragraph{改写坐标与移动点}
同一个计算式可以表示两件不同的事。取$p=(0,\tfrac35,\tfrac45)$，北极图坐标为$(0,3)$。旋转$R_\theta$把$p$送到$(-\tfrac35\sin\theta,\tfrac35\cos\theta,\tfrac45)$，坐标变为$(-3\sin\theta,3\cos\theta)$：点真的移动了，坐标的变化反映了这一移动。坐标变换$\tau=\sigma_S\circ\sigma_N^{-1}$把$(0,3)$变为$(0,\tfrac13)$，点却没有动，只是换了一张图来记录。

用本节的语言说，坐标变换是恒等映射$\mathrm{id}:S^2\to S^2$在输入用$\sigma_N$、输出用$\sigma_S$时的坐标表示。所以判断一个坐标公式描述的是什么，不能只看公式，还要看它是同一映射在不同图下的表示，还是不同映射在同一图下的表示。前者改变的是记录，后者改变的是点；后文讨论微分时，前者只改变切向量的分量，后者把切向量送到另一个点的切空间。
